\begin{lemma}
Let $\Omega\subset\mathbb{C}$ be open, let $H\colon\mathbb{C}\to\mathbb{C}$ be
$C^\infty$ on $\Omega$, and let $i,j\ge0$ be integers.  Then
\begin{enumerate}
\item[(a)] $\dfrac{\partial^{i+j}H}{\partial\xi^i\partial\bar\xi^j}$ is
$C^\infty$ on $\Omega$, and
\item[(b)] $\dfrac{\partial}{\partial\bar\xi}
\dfrac{\partial^{i+j}H}{\partial\xi^i\partial\bar\xi^j}(\xi)
=\dfrac{\partial^{i+j+1}H}{\partial\xi^i\partial\bar\xi^{j+1}}(\xi)$ for every
$\xi\in\Omega$,
\end{enumerate}
where the iterated derivatives are those of \noderef{WirtingerIter} and
$\partial/\partial\bar\xi$ is as in \noderef{DZbar}.
\end{lemma}
\begin{proof}
Throughout, $\partial_x$ and $\partial_y$ denote the two real first partial
derivatives, taken through the real Fr\'echet derivative as in \noderef{DZ};
thus for any $g$ and any point, $\partial_xg=Dg[1]$ and $\partial_yg=Dg[i]$,
and by \noderef{DZ} and \noderef{DZbar}
\begin{equation}\tag{0}
  \frac{\partial g}{\partial\xi}=\tfrac12\bigl(\partial_xg-i\,\partial_yg\bigr),
  \qquad
  \frac{\partial g}{\partial\bar\xi}=\tfrac12\bigl(\partial_xg+i\,\partial_yg\bigr).
\end{equation}

\emph{Step 1: both Wirtinger operators preserve $C^\infty(\Omega)$.}  Let $g$ be
$C^\infty$ on the open set $\Omega$.  Then $g$ is in particular
real-differentiable at each point of $\Omega$, and each of its first partial
derivatives $\partial_xg$, $\partial_yg$ is again $C^\infty$ on $\Omega$ (the
partial derivatives of a $C^\infty$ function on an open set are $C^\infty$
there).  By (0), $\partial g/\partial\xi$ and $\partial g/\partial\bar\xi$ are
fixed $\mathbb{C}$-linear combinations of $\partial_xg$ and $\partial_yg$, hence
both are $C^\infty$ on $\Omega$.

\emph{Step 2: proof of (a).}  By \noderef{WirtingerIter},
$\partial^{i+j}H/\partial\xi^i\partial\bar\xi^j$ is obtained from $H$ by
applying $j$ times the operator $\partial/\partial\bar\xi$ and then $i$ times
the operator $\partial/\partial\xi$.  Each of these $i+j$ operators carries
$C^\infty(\Omega)$ into $C^\infty(\Omega)$ by Step 1, and $H\in C^\infty(\Omega)$
by hypothesis.  Induction on the number of operators applied gives
$\partial^{i+j}H/\partial\xi^i\partial\bar\xi^j\in C^\infty(\Omega)$.

\emph{Step 3: the two operators commute on $C^\infty(\Omega)$.}  Let $g$ be
$C^\infty$ on $\Omega$ and let $\xi\in\Omega$.  Using (0) twice, and writing
$\partial_u\partial_vg$ for $\partial_u(\partial_vg)$,
\begin{align*}
  \frac{\partial}{\partial\xi}\Bigl(\frac{\partial g}{\partial\bar\xi}\Bigr)
  &=\tfrac12\bigl(\partial_x-i\,\partial_y\bigr)
     \Bigl[\tfrac12\bigl(\partial_xg+i\,\partial_yg\bigr)\Bigr]\\
  &=\tfrac14\Bigl(\partial_x\partial_xg+\partial_y\partial_yg
      +i\bigl(\partial_x\partial_yg-\partial_y\partial_xg\bigr)\Bigr),
\end{align*}
and symmetrically
\[
  \frac{\partial}{\partial\bar\xi}\Bigl(\frac{\partial g}{\partial\xi}\Bigr)
  =\tfrac14\Bigl(\partial_x\partial_xg+\partial_y\partial_yg
      -i\bigl(\partial_x\partial_yg-\partial_y\partial_xg\bigr)\Bigr),
\]
both at every point of $\Omega$; in each case the step is the expansion of the
product of the two first-order operators, legitimate because the first
derivatives of $g$ are themselves differentiable on $\Omega$ by Step 1.  Since
$g$ is $C^\infty$ on the open set $\Omega$, its second derivative is symmetric
at every point of $\Omega$ (Schwarz's theorem on mixed partials, valid for
$C^2$ functions on an open set), so
$\partial_x\partial_yg=\partial_y\partial_xg$ on $\Omega$.  The two displayed
expressions therefore coincide:
\begin{equation}\tag{1}
  \frac{\partial}{\partial\xi}\Bigl(\frac{\partial g}{\partial\bar\xi}\Bigr)(\xi)
  =\frac{\partial}{\partial\bar\xi}\Bigl(\frac{\partial g}{\partial\xi}\Bigr)(\xi)
  \qquad\text{for every }\xi\in\Omega .
\end{equation}

\emph{Step 4: proof of (b), by induction on $i$ with $j$ fixed.}

\emph{Base case $i=0$.}  By \noderef{WirtingerIter},
$\partial^{j}H/\partial\xi^0\partial\bar\xi^{j}$ is
$(\partial/\partial\bar\xi)^{j}H$, so applying $\partial/\partial\bar\xi$ to it
gives $(\partial/\partial\bar\xi)^{j+1}H$, which is
$\partial^{j+1}H/\partial\xi^0\partial\bar\xi^{j+1}$.  This is an identity of
functions on all of $\mathbb{C}$ and needs no hypothesis.

\emph{Inductive step.}  Assume
\begin{equation}\tag{2}
  \frac{\partial}{\partial\bar\xi}
  \Bigl(\frac{\partial^{i+j}H}{\partial\xi^i\partial\bar\xi^j}\Bigr)(\zeta)
  =\frac{\partial^{i+j+1}H}{\partial\xi^i\partial\bar\xi^{j+1}}(\zeta)
  \qquad\text{for every }\zeta\in\Omega .
\end{equation}
Put $g=\partial^{i+j}H/\partial\xi^i\partial\bar\xi^j$, which is $C^\infty$ on
$\Omega$ by Step 2.  By \noderef{WirtingerIter},
\[
  \frac{\partial^{i+1+j}H}{\partial\xi^{i+1}\partial\bar\xi^{j}}
  =\frac{\partial g}{\partial\xi},
  \qquad
  \frac{\partial^{i+j+1}H}{\partial\xi^{i}\partial\bar\xi^{j+1}}
  \ \text{ is a $C^\infty(\Omega)$ function by Step 2.}
\]
Let $\xi\in\Omega$.  Then
\[
  \frac{\partial}{\partial\bar\xi}
  \Bigl(\frac{\partial^{i+1+j}H}{\partial\xi^{i+1}\partial\bar\xi^{j}}\Bigr)(\xi)
  =\frac{\partial}{\partial\bar\xi}\Bigl(\frac{\partial g}{\partial\xi}\Bigr)(\xi)
  =\frac{\partial}{\partial\xi}\Bigl(\frac{\partial g}{\partial\bar\xi}\Bigr)(\xi),
\]
by (1) applied to $g$.  By the induction hypothesis (2), the functions
$\partial g/\partial\bar\xi$ and
$\partial^{i+j+1}H/\partial\xi^{i}\partial\bar\xi^{j+1}$ agree at every point
of $\Omega$.  Since $\Omega$ is open and $\xi\in\Omega$, two functions agreeing
on $\Omega$ agree on a neighbourhood of $\xi$, hence have the same real
Fr\'echet derivative at $\xi$ and therefore, by (0), the same
$\partial/\partial\xi$ at $\xi$.  Consequently
\[
  \frac{\partial}{\partial\xi}\Bigl(\frac{\partial g}{\partial\bar\xi}\Bigr)(\xi)
  =\frac{\partial}{\partial\xi}
   \Bigl(\frac{\partial^{i+j+1}H}{\partial\xi^{i}\partial\bar\xi^{j+1}}\Bigr)(\xi)
  =\frac{\partial^{i+1+j+1}H}{\partial\xi^{i+1}\partial\bar\xi^{j+1}}(\xi),
\]
the last equality again by \noderef{WirtingerIter}.  Combining the two displays
gives (b) for $i+1$, and the induction is complete.
\end{proof}
